
\section{本轮新结果：正部的行方差机制、天花板 $L(N)$，与旧队首 P1 的否证}\label{sec:rowvar}

本节是第 12 轮工作的完整记录。它有一条\textbf{坏消息}（旧队首 P1 的「腔掩码不存在性」路线被一个精确有理证书否证，
第~\ref{sec:chamberdead} 小节），和一条\textbf{好消息}（我们找到了一个全新的不等式机制——「行方差机制」——
它同时把 Mohar 的 $\frac{1+\sqrt3}2$ 与 \STW\ 的 $\Phi=\frac{3+\sqrt5}4$ 重导为\emph{同一个}闭式 $L(N)$ 在
$N\to\infty$ 与 $N=6$ 处的取值，并且在冠军处\emph{四重取等}；第~\ref{sec:rowid}--\ref{sec:symlayer} 小节）。
坏消息的意义在于：它说明「纯几何的存在性/不存在性」这条路上已经没有任何可关闭的类；
好消息的意义在于：它把主线缺口从「一条越来越难的不等式」换成了「一个可以逐条加约束的有限维松弛」，
并且\textbf{把缺口定量地压缩到 $3.2\times10^{-3}$}（在 $Z_5$-对称层内）。

\subsection{旧队首 P1 被否证：$K_{2,2,2}$ 是腔掩码（精确有理证书）}\label{sec:chamberdead}

注~\ref{rmk:tab19read}(iv) 把「$\Sigma\notin\mathcal S_6^{\mathrm{int}}$」\emph{等价}化为
\eqref{eq:chamber6}（$\Sigma\ne1\{Q\ge0\}$ 对一切 $Q\in\mathcal P_3(6)$）。这个等价化是\textbf{错的}。

\begin{proposition}[腔掩码否证，已证（$\Q$ 上精确）]\label{prop:chamberdead}
设 $M$ 为 $K_6$ 中某个完美匹配的置换矩阵（三对不相邻的 $2$-循环），令
\[
Q=\tfrac12\,(I_6-M)\ \in\ \Q^{6\times6}.
\]
则 $Q^\top=Q$、$Q^2=Q$、$\tr Q=3$（故 $Q\in\mathcal P_3(6)$，特征值为 $1,1,1,0,0,0$），
$q_{ii}=\frac12>0$，匹配对上 $q_{ij}=-\frac12<0$，其余位置 $q_{ij}=0$。于是
\[
1\{Q\ge0\}=J_6-M=I_6+A(K_{2,2,2}),
\]
即 \texttt{EznW}$\,=K_{2,2,2}$ \textbf{是}腔掩码。因此 \eqref{eq:chamber6} 对 $\Sigma=I+A(K_{2,2,2})$ 为\textbf{假}。
\end{proposition}

\begin{proof}
$M$ 是对合（$M^2=I_6$）且 $\tr M=0$，特征值为 $+1$（三重）与 $-1$（三重）。
故 $Q=\frac12(I_6-M)$ 的特征值为 $\frac12(1-1)=0$（三重）与 $\frac12(1+1)=1$（三重），
$Q^2=\frac14(I-2M+M^2)=\frac14(2I-2M)=Q$，$\tr Q=3$。符号模式由构造直接读出；
$K_6$ 去掉一个完美匹配正是完全三部图 $K_{2,2,2}$。全部计算在 $\Q$ 上进行，无任何数值成分。$\square$
\end{proof}

\begin{remark}[这条否证的真实后果]\label{rmk:chamberdead}
(i) 我们把 $N=6$ 的 $19$ 个未闭合类逐个做了数值腔搜索，结果是
\textbf{$19/19$ 都是腔掩码}（每个类都找到 $Q\in\mathcal P_3(6)$ 使 $1\{Q\ge0\}=\Sigma$，
符号违背量 $\le5\times10^{-19}$，即精确到机器精度）。所以
\textbf{「证明不是腔掩码」这条几何路线关闭的类数是 $0$}，旧 P1 应当弹出。

(ii) 引理~\ref{lem:bdryrig} 本身没有错，错的是对它的读法：对无支配点的类，
「$\Sigma$ 是腔掩码」只是「$\Sigma\in\mathcal S_6^{\mathrm{int}}$」的\textbf{必要条件}，不是充分条件。
真正缺的约束是\emph{耦合}条件
\[
\Sigma=1\bigl\{\Pi_3(D_p^{1/2}\Sigma D_p^{1/2})\ge0\bigr\},
\]
即掩码必须是\textbf{它自己}（带某个正权 $p$）的前三特征投影的符号模式。
命题~\ref{prop:chamberdead} 的 $Q=\frac12(I_6-M)$ 恰好说明了二者之差：
它的 $1\{Q\ge0\}=I+A(K_{2,2,2})$，但 $Q_+=\frac12I_6$ 的 Perron 向量是退化的、
$\rho(Q_+)=\frac12$，与 $F(I+A(K_{2,2,2}))=\kappa_3(5)=\frac{9+4\sqrt6}{15}$ 毫无关系。
「是腔掩码」完全不涉及自洽性，这就是它太弱的原因。

(iii) 因此注~\ref{rmk:tab19read}(iv) 中的「等价」必须降级为「必要条件」；
\eqref{eq:s6int} 仍然是正确的蕴含（$\Leftarrow$ 方向不受影响），但它的\emph{后件}不再能用
引理~\ref{lem:chamber} 的纯几何陈述来建立。$N=6$ 的 $19$ 类回到「必须使用自洽耦合」的状态。
\end{remark}

\subsection{本节记号与三条新恒等式}\label{sec:rowid}

\textbf{记号（本节专用，与 \S\ref{sec:nogo} 的框架行记号分离）。}
设 $Q=(q_{ij})\in\mathcal P_3(N)$，即 $Q^2=Q=Q^\top$、$\tr Q=3$（故 $0\preceq Q\preceq I$），
$x\ge0$、$\|x\|=1$、$Q_+x=\rho x$（$\rho=\rho(Q_+)=\lambda_{\max}(Q_+)$）。
沿用命题~\ref{prop:perronrig} 的 $\alpha=x^\top Qx$、$\beta=x^\top Q_-x$、$d=Q_-x$、$\delta=\|d\|^2$、$\ell=x^\top|Q|x$。
新增：
\[
g_i:=q_{ii}=c_i^2\quad(\text{故 }0\le g_i\le1,\ \textstyle\sum_ig_i=3),\qquad
\pi_i:=x_i^2\quad(\textstyle\sum_i\pi_i=1),
\]
\[
O:=\{(i,j):i\ne j,\ q_{ij}<0\},\qquad A:=\{(i,j):i\ne j,\ q_{ij}>0\},
\]
\[
a_i:=\!\!\sum_{j\ne i}\!\bigl((q_{ij})_+\bigr)^2,\quad
r_i:=\!\!\sum_{j\ne i}\!\bigl((q_{ij})_-\bigr)^2,\quad
v_i:=\!\!\sum_{(i,j)\in A}\!\!\pi_j,\quad
w_i:=\!\!\sum_{(i,j)\in O}\!\!\pi_j,
\]
于是 $v_i+w_i+\pi_i=1$（$\pi_i$ 自己那一项既不属于 $A$ 也不属于 $O$）。再记
\[
s_i:=\sum_{j\ne i}(q_{ij})_+x_j\quad(\text{正部行和}),\qquad
S_O:=\sum_ir_i,\qquad W_O:=\sum_i\pi_iw_i,\qquad R:=\sum_ir_iw_i .
\]

\begin{lemma}[三条新恒等式，已证]\label{lem:rowid}
在上述记号下，对\emph{任意} $Q\in\mathcal P_3(N)$ 与\emph{任意}满足 $Q_+x=\rho x$ 的 $x\ge0$、$\|x\|=1$：
\begin{enumerate}[label=(\roman*),leftmargin=2.2em]
\item \textbf{行能量恒等式}：$\displaystyle a_i+r_i=\sum_{j\ne i}q_{ij}^2=g_i(1-g_i)$；
特别地 $\sum_{ij}q_{ij}^2=3$、$\sum_i(a_i+r_i)=3-\sigma$，其中 $\sigma:=\sum_ig_i^2$。
\item \textbf{Perron 行方程}：$s_i=(\rho-g_i)x_i$；等价地 $\rho=g_i+s_i/x_i$（$x_i>0$ 时）。
\item \textbf{Hadamard 平方与三次迹}：$Q\circ Q\succeq0$、$(Q\circ Q)\one=g$、$\one^\top(Q\circ Q)\one=3$，
故 $\lambda_{\max}(Q\circ Q)\le\max_ig_i\le1$；又 $\tr Q^3=\tr Q=3$ 给出\textbf{三角形恒等式}
\begin{equation}\label{eq:tri}
\sum_{i<j<k}q_{ij}q_{jk}q_{ki}\ =\ \tfrac16\Bigl(3-3\sigma+2\sum_ig_i^3\Bigr).
\end{equation}
\end{enumerate}
\end{lemma}

\begin{proof}
(i) $(Q^2)_{ii}=\sum_jq_{ij}^2=Q_{ii}=g_i$，而 $q_{ii}^2=g_i^2$，故 $\sum_{j\ne i}q_{ij}^2=g_i-g_i^2$；
又 $q_{ij}^2=((q_{ij})_+)^2+((q_{ij})_-)^2$ 逐项成立。对 $i$ 求和并用 $\sum g_i=3$ 得第二句。

(ii) $Q_+x=\rho x$ 的第 $i$ 行是 $(Q_+)_{ii}x_i+\sum_{j\ne i}(q_{ij})_+x_j=\rho x_i$；
因 $g_i\ge0$ 有 $(Q_+)_{ii}=g_i$，故 $s_i=(\rho-g_i)x_i$。

(iii) $Q\succeq0\Rightarrow Q\circ Q\succeq0$（Schur 乘积定理）。
$(Q\circ Q)\one$ 的第 $i$ 分量就是 $\sum_jq_{ij}^2=g_i$，即 (i)；
再与 $\one$ 作内积得 $\one^\top(Q\circ Q)\one=\sum_ig_i=3$。
非负矩阵的谱半径介于最小行和与最大行和之间，故 $\lambda_{\max}(Q\circ Q)\le\max_ig_i$。
最后，$\tr Q^3=\sum_{ijk}q_{ij}q_{jk}q_{ki}$ 按「下标重合个数」分三类：
三个都相同给 $\sum_ig_i^3$；恰有两个相同给 $3\sum_{i\ne j}g_iq_{ij}^2=3\sum_ig_i\cdot g_i(1-g_i)=3\sigma-3\sum_ig_i^3$（用 (i)）；
三个互不相同给 $6\sum_{i<j<k}q_{ij}q_{jk}q_{ki}$。三者之和 $=\tr Q^3=\tr Q=3$，整理即 \eqref{eq:tri}。$\square$
\end{proof}

\begin{remark}\label{rmk:rowidread}
(i) 引理~\ref{lem:rowid}(i) 是\textbf{秩三投影的全部「逐行」内容}：它说每一行的能量 $g_i(1-g_i)$
被正部与负部按 $a_i:r_i$ 精确瓜分。第~\ref{sec:stwanat} 节的 \STW\ 解剖只用到了它的求和版
$\sum_{ij}q_{ij}^2=3$（即 Mohar 界里的那个 $\sqrt3$）；\textbf{逐行}使用它是本轮的新想法。

(ii) (ii) 把「$\rho$ 有多大」直接翻译成「每一行的正部行和有多大」，
于是对正部行和做一次 Cauchy--Schwarz 就得到 $\rho$ 的\emph{逐行}上界——这是第~\ref{sec:relax} 小节的发动机。

(iii) \eqref{eq:tri} 是一条\textbf{尚未使用}的强约束：它说全部三角形乘积之和由对角线单独决定。
冠军处 $\sigma=1.533964$、$\sum g_i^3=0.796424$，故右端 $=-0.0015073$（一个很小的负数），
这说明钝角三角形与锐角三角形几乎完全抵消。把它翻译成符号模式的组合约束是 P1$'$ 的入口之一。
\end{remark}
